\chapter{Getting Started: Insertion Sort and Loop Invariants}

There is a reason almost every serious book on algorithms begins with sorting,
and a reason this one begins with \emph{insertion sort} in particular. Sorting is
the cleanest possible playground. The problem is easy to state, the answer is easy
to check, and yet the act of sorting forces us to confront, in miniature, almost
every idea that will haunt us for the rest of the book: how to describe a process
precisely, how to convince ourselves it is \emph{correct}, and how to measure how
\emph{long} it takes without ever running it on a computer.

Insertion sort is the algorithm you already use, whether you know it or not, when
you pick up a hand of playing cards and tidy them. That is not a coincidence. We
are going to take that everyday motion, slow it down to a crawl, and turn it into
something we can reason about with full mathematical confidence. By the end of this
chapter you will have a tool --- the \emph{loop invariant} --- that lets you prove a
loop does what you claim, and a habit --- counting \emph{order of growth} --- that lets
you compare algorithms before you have written a single line of runnable code.

\section{The problem we are solving}

Let us be painfully precise about the goal, because precision here pays off
everywhere later. We are given a list of numbers, and we want them rearranged into
increasing order. ``Increasing'' will mean ``non-decreasing'': equal values are
allowed to sit next to each other.

\begin{formalbox}
\textbf{The sorting problem.}
\emph{Input:} a sequence of $n$ numbers $\langle a_1, a_2, \dots, a_n\rangle$.

\emph{Output:} a permutation (a reordering) $\langle a'_1, a'_2, \dots, a'_n\rangle$
of the input such that
\[
a'_1 \le a'_2 \le \cdots \le a'_n .
\]
\end{formalbox}

Two words in that box are doing quiet but heavy lifting, and you should not let them
slide past you.

\begin{focusbox}
A correct sorting algorithm must satisfy \emph{two} conditions, not one.
\begin{enumerate}
  \item \textbf{Ordered:} the output is sorted, $a'_1 \le \cdots \le a'_n$.
  \item \textbf{Permutation:} the output contains exactly the same elements as the
        input, just possibly in a different order --- nothing added, nothing lost,
        nothing duplicated.
\end{enumerate}
An algorithm that returns $\langle 1,2,3,\dots,n\rangle$ no matter what you feed it
is perfectly ordered and completely wrong. Hold on to \emph{both} conditions.
\end{focusbox}

We will treat the numbers as living in an array, which is just a row of labeled
boxes. We write $A[1\,..\,n]$ for an array with $n$ slots, and $A[i]$ for the value in
slot $i$. Throughout this book array indices start at $1$, matching the way you would
count cards in your hand and matching CLRS; later chapters will note where real
programming languages start at $0$ instead.

\begin{historybox}
This chapter mirrors \textbf{CLRS, Chapter 2 (``Getting Started'')}. CLRS opens with
exactly insertion sort, the loop-invariant method, running-time analysis, and merge
sort as the first taste of divide-and-conquer. We will cover the same ground but
linger far longer on the \emph{why}. Sorting itself is ancient as a human activity
(librarians, clerks, and card-sorters did it for centuries); the formal study of
sorting algorithms blossomed in the 1940s and 1950s alongside the first stored-program
computers, and Donald Knuth's \emph{The Art of Computer Programming, Volume 3} (1973)
remains the encyclopedic reference.
\end{historybox}

\section{The card-player's idea}

Picture yourself dealt a hand of cards, face down, on the table. You pick them up one
at a time with your right hand and slot each new card into its correct place among the
cards already arranged in your left hand. At every moment, the cards in your left hand
are sorted; the cards still on the table are not yet looked at.

\begin{intuitionbox}
Here is the whole algorithm in one sentence: \emph{keep a sorted region on the left,
and repeatedly take the next unsorted item and slide it leftward, past everything
bigger than it, until it lands in the right spot.}

The sorted region starts as just the first card (a single card is trivially sorted).
Each step grows the sorted region by exactly one. After you have inserted the last
card, the whole hand is sorted. That ``grow the sorted region by one each time'' is
the heartbeat of the algorithm, and it is exactly what we will pin down with a loop
invariant.
\end{intuitionbox}

The only slightly clever part is \emph{how} a card slides into place. You do not yank
out a gap and drop it in. Instead you hold the new card (call it the \emph{key}),
compare it with the card to its left, and if that left card is bigger, you shove the
left card one slot to the right to make room, then keep comparing with the next card
further left. You stop when you meet a card that is not bigger than the key (or you run
off the left end), and you drop the key into the gap you just opened.

Figure~\ref{ch02:fig-hand} shows the two regions at one frozen instant: a sorted left
hand, an untouched pile on the right, and the single key being lifted across the
boundary between them.

\begin{figure}[H]
\centering
\begin{tikzpicture}[
    box/.style={draw, minimum size=8mm, font=\small},
    sorted/.style={box, fill=targetgreen!18},
    unsorted/.style={box, fill=gray!10},
    key/.style={box, fill=quotientorange!30},
]
\node[sorted] (s1) at (0,0) {2};
\node[sorted] (s2) at (0.85,0) {4};
\node[sorted] (s3) at (1.7,0) {7};
\node[key]    (k)  at (3.0,0.9) {5};
\node[unsorted](u1) at (3.4,0) {1};
\node[unsorted](u2) at (4.25,0) {9};
\draw[thick,dashed] (2.55,-0.6) -- (2.55,0.6);
\draw[->,thick,quotientorange!75!black] (u1.north) .. controls (3.2,0.55) .. (k.south east);
\node[font=\footnotesize,targetgreen!50!black] at (0.85,-0.85) {sorted left hand};
\node[font=\footnotesize,gray] at (3.825,-0.85) {untouched pile};
\node[font=\footnotesize,quotientorange!60!black,align=center] at (3.0,1.65) {key lifted,\\ seeking its slot};
\end{tikzpicture}
\caption{One frozen instant of insertion sort. The dashed line is the boundary; the key
just left the unsorted pile and will slide left until it meets a card no bigger than
itself.}
\label{ch02:fig-hand}
\end{figure}

\section{Insertion sort, written down}

Let us turn that motion into precise pseudocode. We process positions $j = 2, 3,
\dots, n$. At the start of step $j$, the subarray $A[1\,..\,j-1]$ is already sorted (the
cards in your left hand), and $A[j]$ is the new card we must insert.

\begin{algorithm}[H]
\caption{\textsc{Insertion-Sort}$(A)$}
\KwIn{An array $A[1\,..\,n]$ of $n$ numbers.}
\KwOut{The same array, rearranged so that $A[1] \le A[2] \le \cdots \le A[n]$.}
\For{$j \leftarrow 2$ \KwTo $n$}{
  $\mathit{key} \leftarrow A[j]$\;
  \tcp{Insert $A[j]$ into the sorted sequence $A[1\,..\,j-1]$.}
  $i \leftarrow j - 1$\;
  \While{$i > 0$ \textbf{and} $A[i] > \mathit{key}$}{
    $A[i+1] \leftarrow A[i]$\;
    $i \leftarrow i - 1$\;
  }
  $A[i+1] \leftarrow \mathit{key}$\;
}
\Return $A$\;
\end{algorithm}

Read it slowly. The outer \textbf{for} loop walks the boundary of the sorted region
rightward. Inside, $\mathit{key}$ is the card we are placing. The \textbf{while} loop
does the shoving: as long as we have not fallen off the left edge ($i > 0$) and the
card at $A[i]$ is strictly bigger than the key, we copy $A[i]$ one slot to the right
(opening room) and step $i$ leftward. When the loop stops, slot $i+1$ is the empty gap,
and that is exactly where the key belongs, so we drop it there.

\begin{trapbox}
Notice the comparison is $A[i] > \mathit{key}$, a \emph{strict} inequality. This matters
for two reasons. First, it makes the sort \emph{stable}: when two values are equal, the
one that started earlier stays earlier, because we stop shoving the moment we meet an
equal value rather than sliding past it. Second, if you accidentally wrote $A[i] \ge
\mathit{key}$, the sort would still be \emph{correct} (still sorted, still a
permutation) but no longer stable, and on a run of equal keys it would do needless
extra work. Stability will matter a great deal in the later chapter on sorting in
linear time, where we lean on it to chain sorting passes together.
\end{trapbox}

\section{A full numeric trace}

Abstractions feel weightless until you push real numbers through them, so let us sort
\[
A = \langle 5,\; 2,\; 4,\; 6,\; 1,\; 3 \rangle .
\]
We will show the array after each pass of the outer loop. The vertical bar
``$\mid$'' marks the boundary: everything to its left is the sorted region, everything
to its right is untouched.

\begin{labbox}
\textbf{Trace of \textsc{Insertion-Sort} on $\langle 5,2,4,6,1,3\rangle$.}

Start: $\;5 \mid 2\;\;4\;\;6\;\;1\;\;3$ \hfill (sorted region is just $A[1]=5$)

\smallskip
\textbf{$j=2$, key $=2$.} Compare with $5$: since $5>2$, shove $5$ right. Now $i=0$,
run off the edge, drop $2$ in slot $1$.
\[
2\;\;5 \mid 4\;\;6\;\;1\;\;3
\]

\textbf{$j=3$, key $=4$.} Compare with $5$: $5>4$, shove $5$ right. Compare with $2$:
$2 \not> 4$, stop. Drop $4$ after the $2$.
\[
2\;\;4\;\;5 \mid 6\;\;1\;\;3
\]

\textbf{$j=4$, key $=6$.} Compare with $5$: $5 \not> 6$, stop immediately. The key is
already in place; drop it back where it was.
\[
2\;\;4\;\;5\;\;6 \mid 1\;\;3
\]

\textbf{$j=5$, key $=1$.} Compare with $6,5,4,2$ in turn; each is bigger, so all four
shove right. Run off the edge, drop $1$ in slot $1$.
\[
1\;\;2\;\;4\;\;5\;\;6 \mid 3
\]

\textbf{$j=6$, key $=3$.} Compare with $6,5,4$ (all bigger, shove right), then $2$:
$2 \not> 3$, stop. Drop $3$ after the $2$.
\[
1\;\;2\;\;3\;\;4\;\;5\;\;6 \mid
\]
\end{labbox}

The sorted region grew by one on every pass --- $1,2,3,4,5,6$ elements --- and after
the last pass the whole array is sorted. Notice two passes worth contrasting: at $j=4$
the key $6$ caused \emph{zero} shoves (it was already in place), while at $j=5$ the key
$1$ caused \emph{four} shoves (it had to travel all the way to the front). That gap
between ``barely any work'' and ``maximum work'' is precisely what the running-time
analysis later in the chapter will measure.

The hardest pass of this trace is $j=5$, where the key $1$ travels all the way to the
front. Figure~\ref{ch02:fig-slide} captures that pass: the four boxed values $6,5,4,2$
each get shoved one step right, and $1$ drops into the slot that opens up at the front.

\begin{figure}[H]
\centering
\begin{tikzpicture}[
    box/.style={draw, minimum size=8mm, font=\small},
    sorted/.style={box, fill=targetgreen!16},
    unsorted/.style={box, fill=gray!8},
    key/.style={box, fill=quotientorange!30},
]
\node[align=right,font=\small] at (-2.7,0.2) {After $j{=}4$, inserting\\ $\mathit{key}=A[5]=1$:};
\node[sorted] (a1) at (0,0) {2};
\node[sorted] (a2) at (0.85,0) {4};
\node[sorted] (a3) at (1.7,0) {5};
\node[sorted] (a4) at (2.55,0) {6};
\node[key]    (a5) at (3.4,0) {1};
\node[unsorted](a6) at (4.25,0) {3};
\draw[->,thick,quotientorange!70!black] (a5.north) .. controls (2.1,1.0) and (-0.6,1.0) .. (a1.north west);
\node[font=\small] at (1.4,1.2) {key slides left past everything bigger};
\node[font=\footnotesize] at (0,-0.65) {1};
\node[font=\footnotesize] at (0.85,-0.65) {2};
\node[font=\footnotesize] at (1.7,-0.65) {3};
\node[font=\footnotesize] at (2.55,-0.65) {4};
\node[font=\footnotesize] at (3.4,-0.65) {5};
\node[font=\footnotesize] at (4.25,-0.65) {6};
\node[font=\footnotesize,gray] at (5.6,-0.65) {(index)};
\end{tikzpicture}
\caption{The $j=5$ pass on $\langle 5,2,4,6,1,3\rangle$. The key $1$ is smaller than all
four sorted values, so each is shifted one slot right and $1$ lands in slot $1$ --- four
shoves, the most this pass can ever cost.}
\label{ch02:fig-slide}
\end{figure}

\section{Loop invariants: how we \emph{know} it works}

Tracing one example builds confidence, but it does not build certainty. The array
$\langle 5,2,4,6,1,3\rangle$ came out right --- but what about the millions of arrays we
did not try? We need an argument that covers \emph{all} inputs at once. The standard
tool for reasoning about loops is the \emph{loop invariant}: a statement that is true
every time the loop is about to start a new iteration, no matter what the input was.

\begin{intuitionbox}
A loop invariant is like a promise you make about the state of the world at the top of
the loop, and then keep, iteration after iteration. If the promise is true before the
first iteration, and each iteration preserves it, then by the time the loop ends the
promise still holds --- and we engineer the promise so that, combined with the reason
the loop stopped, it forces the answer we want. It is mathematical induction wearing
work clothes.
\end{intuitionbox}

To use a loop invariant we check three things. The names are worth memorizing because
you will perform this ritual for every loop in the book.

\begin{focusbox}
\textbf{The three obligations of a loop invariant.}
\begin{itemize}
  \item \textbf{Initialization.} It is true before the first iteration of the loop.
        (This is the base case of the induction.)
  \item \textbf{Maintenance.} If it is true before an iteration, it remains true before
        the next iteration. (This is the inductive step.)
  \item \textbf{Termination.} When the loop ends, the invariant --- together with the
        condition that ended the loop --- gives us a useful property that proves the
        algorithm correct.
\end{itemize}
\end{focusbox}

The first two are exactly induction: a base case and a step. The third is the twist
that distinguishes a loop invariant from plain induction --- we deliberately choose the
invariant so that its \emph{ending} state tells us what we set out to prove. An
ordinary induction proves a fact for all $n$; a loop invariant additionally cashes that
fact out into a statement about the finished computation.

\subsection{The invariant for insertion sort}

Here is the invariant for the outer \textbf{for} loop of \textsc{Insertion-Sort}.

\begin{formalbox}
\textbf{Outer-loop invariant.} At the start of each iteration of the \textbf{for} loop
(the one indexed by $j$), the subarray $A[1\,..\,j-1]$ consists of the elements
originally in $A[1\,..\,j-1]$, but in sorted order.
\end{formalbox}

Read it twice. It claims two things --- sorted, \emph{and} the same elements --- which
echoes our two correctness conditions from the very start. Now we discharge the three
obligations.

\begin{formalbox}
\textbf{Initialization.} Before the first iteration we have $j = 2$, so the subarray in
question is $A[1\,..\,1]$, a single element. A list of one element is trivially in
sorted order, and it obviously contains the original element $A[1]$. The invariant
holds.

\textbf{Maintenance.} Assume the invariant holds at the start of an iteration:
$A[1\,..\,j-1]$ is a sorted rearrangement of the original $A[1\,..\,j-1]$. The body
grabs $\mathit{key} = A[j]$ and runs the \textbf{while} loop. That inner loop slides
each element of $A[1\,..\,j-1]$ that is strictly greater than $\mathit{key}$ one
position to the right, and then deposits $\mathit{key}$ into the freed slot. Because
every element moved was greater than $\mathit{key}$, and every element to the left of
the landing slot was not greater than $\mathit{key}$, the result $A[1\,..\,j]$ is
sorted; and since we only moved elements and re-inserted the very value we removed, it
contains exactly the original elements of $A[1\,..\,j]$. Incrementing $j$ for the next
iteration restores the invariant's form ``$A[1\,..\,j-1]$ is a sorted rearrangement''
--- now over a region one longer.

\textbf{Termination.} The loop ends when $j$ exceeds $n$, i.e. when $j = n+1$.
Substituting $j = n+1$ into the invariant, the subarray $A[1\,..\,n]$ is a sorted
rearrangement of the original elements. But $A[1\,..\,n]$ is the entire array. So the
whole array is sorted and is a permutation of the input --- which is exactly the sorting
problem. $\qed$
\end{formalbox}

That is a complete correctness proof. Pause to appreciate what just happened: three
short paragraphs covered every possible input array of every possible length. No
amount of testing could ever do that.

\begin{trapbox}
A subtlety swept under the rug above: maintenance leaned on the inner \textbf{while}
loop ``doing the right thing.'' To be fully rigorous you would give the \textbf{while}
loop its \emph{own} invariant --- for example, ``$A[i+1\,..\,j]$ all hold values $\ge
\mathit{key}$, shifted from their original spots, and $A[1\,..\,i]$ is unchanged and
still sorted.'' We will not write that one out in full here (it is an exercise), but you
should notice that \emph{nested} loops generally need \emph{nested} invariants. Skipping
the inner argument is a common way a ``proof'' quietly springs a leak.
\end{trapbox}

\begin{walkbox}
Take this on a walk. We claimed a single-element list is ``trivially sorted.'' What
about an \emph{empty} list --- is it sorted? Our loop never starts with an empty sorted
region (it starts at $j=2$ with one element). But many algorithms do. Convince yourself
that the statement ``$a'_1 \le a'_2 \le \cdots \le a'_n$'' is \emph{vacuously true} when
$n = 0$ or $n = 1$, because there are no adjacent pairs to violate it. Why is ``vacuously
true'' the right default rather than a cop-out?
\end{walkbox}

\section{How long does it take?}

We have an algorithm we trust. The next question is how \emph{expensive} it is. We could
time it on a laptop, but that answer would depend on the laptop, the programming
language, the weather in the data center, and a hundred other things we do not care
about. We want a measure of cost that depends only on the algorithm and the size of the
input. The classic move is to count \emph{primitive steps}: each line of pseudocode is
assumed to take some constant amount of time per execution, and we add up how many
times each line runs.

Let $n$ be the array length. Let $t_j$ be the number of times the \textbf{while}-loop
test is evaluated for that value of $j$ (this is where all the input-dependence hides
--- some keys slide far, some not at all). Annotating each line with a per-execution
cost $c_k$ and an execution count gives the following.

\begin{center}
\renewcommand{\arraystretch}{1.35}
\begin{tabular}{l l l}
\toprule
\textbf{Line} & \textbf{cost} & \textbf{times} \\
\midrule
\texttt{for $j \leftarrow 2$ to $n$}      & $c_1$ & $n$ \\
\quad $\mathit{key} \leftarrow A[j]$        & $c_2$ & $n-1$ \\
\quad $i \leftarrow j-1$                    & $c_3$ & $n-1$ \\
\quad \texttt{while $i>0$ and $A[i]>key$}   & $c_4$ & $\sum_{j=2}^{n} t_j$ \\
\quad\quad $A[i+1] \leftarrow A[i]$         & $c_5$ & $\sum_{j=2}^{n} (t_j-1)$ \\
\quad\quad $i \leftarrow i-1$               & $c_6$ & $\sum_{j=2}^{n} (t_j-1)$ \\
\quad $A[i+1] \leftarrow \mathit{key}$      & $c_7$ & $n-1$ \\
\bottomrule
\end{tabular}
\end{center}

The \textbf{for} test runs $n$ times (the extra one is the final test that fails and
exits). The body lines run $n-1$ times. The \textbf{while} test for a given $j$ runs
$t_j$ times; the two lines \emph{inside} the \textbf{while} run one fewer, $t_j - 1$
times, because the last test is the one that fails. Adding it all up, the total running
time $T(n)$ is
\[
T(n) = c_1 n + (c_2 + c_3 + c_7)(n-1) + c_4 \sum_{j=2}^{n} t_j
       + (c_5 + c_6) \sum_{j=2}^{n} (t_j - 1).
\]
Everything except the $t_j$ sums is fixed by $n$ alone. The drama lives entirely in
``how far did each key have to slide,'' i.e. in the $t_j$.

\subsection{Best case: already sorted}

If the input is already sorted, then for each $j$ the key $A[j]$ is $\ge$ everything to
its left, so the \textbf{while} test fails immediately: $A[i] = A[j-1] \not> A[j]$.
That means $t_j = 1$ for every $j$. The two sliding lines never run. Substituting,
\[
T(n) = c_1 n + (c_2 + c_3 + c_4 + c_7)(n-1) = an + b
\]
for constants $a, b$. This is a \emph{linear} function of $n$. The best case is fast:
insertion sort barely glances at an already-sorted array.

\subsection{Worst case: reverse-sorted}

Now suppose the input is in strictly \emph{decreasing} order. Every new key is smaller
than everything to its left, so it must slide all the way to the front. For position
$j$ the \textbf{while} test runs $t_j = j$ times (it compares against all $j-1$ earlier
elements and then once more to fall off the edge). We need two standard sums; the
arithmetic series shows up constantly, so meet it now:
\[
\sum_{j=2}^{n} j = \left(\sum_{j=1}^{n} j\right) - 1 = \frac{n(n+1)}{2} - 1,
\qquad
\sum_{j=2}^{n} (j-1) = \sum_{k=1}^{n-1} k = \frac{n(n-1)}{2}.
\]
Plugging these in, every term that involved $\sum t_j$ now contributes something
proportional to $n^2$, and after collecting constants we get
\[
T(n) = a n^2 + b n + c
\]
for some constants $a > 0$, $b$, $c$. This is a \emph{quadratic} function of $n$. The
worst case is slow: doubling the input roughly \emph{quadruples} the work.

\begin{intuitionbox}
Why $n^2$, intuitively? In the worst case, inserting the $j$-th key costs about $j$
comparisons, and we sum that over all $j$ from $2$ to $n$. Summing $1 + 2 + 3 + \cdots +
n$ gives about $n^2/2$. The ``triangle'' of work --- a little for the early keys, a lot
for the late ones --- has area proportional to $n^2$. Any algorithm whose $k$-th step
costs about $k$ will show this same triangular, quadratic behavior.
\end{intuitionbox}

Figure~\ref{ch02:fig-triangle} makes the triangle literal: bar $j$ has height $t_j = j$,
and the bars together fill a right triangle of area about $\tfrac12 n^2$.

\begin{figure}[H]
\centering
\begin{tikzpicture}[scale=0.58]
\foreach \j in {1,...,9} {
  \fill[sourceblue!25] (\j,0) rectangle (\j+0.8,\j);
  \draw[sourceblue!70!black] (\j,0) rectangle (\j+0.8,\j);
}
\draw[->] (0.6,0) -- (11,0) node[right] {\small key index $j$};
\draw[->] (0.6,0) -- (0.6,10) node[above] {\small work $t_j$};
\draw[densely dashed,trapred] (1,1) -- (9.8,9.8);
\node[trapred] at (7.4,4.0) {\small area $\approx \tfrac{1}{2} n^2$};
\end{tikzpicture}
\caption{Worst-case work of insertion sort. The $j$-th key costs about $j$ comparisons,
so the bars climb linearly and their total area --- the running time --- grows like
$\tfrac12 n^2$, hence $\Theta(n^2)$.}
\label{ch02:fig-triangle}
\end{figure}

\subsection{Average case}

What about a ``typical'' random input? On average, when we insert $A[j]$ we expect it
to land somewhere in the middle of the $j-1$ already-sorted elements, so we expect to
shove about half of them: $t_j \approx j/2$. Halving the worst-case sum still leaves a
sum proportional to $n^2$. So the average case is \emph{also} quadratic --- about half
the worst-case constant, but the same shape. This is a recurring and slightly
disheartening lesson:

\begin{focusbox}
For insertion sort, best case is linear ($\Theta(n)$) but \emph{both} the worst case
\emph{and} the average case are quadratic ($\Theta(n^2)$). Being lucky with the average
does not save you here --- ``half of a lot'' is still a lot. To beat $n^2$ on typical
input we will need a genuinely different idea, which is where merge sort comes in.
\end{focusbox}

\section{Order of growth}

Look back at the best- and worst-case formulas: $an + b$ and $an^2 + bn + c$. When $n$
is large, the leading term swamps everything else. For $n = 1000$, the term $n^2 =
1{,}000{,}000$ dwarfs any reasonable $bn$ or $c$; and the constant $a$ in front, while
it affects the stopwatch, does not change the fundamental \emph{shape} of the growth.
So we throw away the lower-order terms and the leading constant, and we keep only the
shape. We say insertion sort's worst-case running time is \emph{of order} $n^2$, and we
will write this $\Theta(n^2)$ (``theta of $n$ squared'').

\begin{intuitionbox}
Order of growth answers the only question that survives changing computers: ``if I make
the input twice as big, roughly how much more work is it?'' For a $\Theta(n)$ algorithm,
double the input means double the work. For $\Theta(n^2)$, double the input means
\emph{four times} the work. For $\Theta(\log n)$, double the input means just \emph{one
more} step. The constants and small terms decide whether you wait $2$ seconds or $5$;
the order of growth decides whether the algorithm is usable at all when $n$ gets large.
\end{intuitionbox}

We will make $\Theta$, $O$, and $\Omega$ fully precise in the next chapter, with
definitions and proofs. For now, carry this loose-but-correct intuition:
\begin{itemize}
  \item $\Theta(f(n))$ means ``grows like $f(n)$'' --- a tight, two-sided bound.
  \item $O(f(n))$ means ``grows no faster than $f(n)$'' --- an upper bound, a ceiling.
  \item $\Omega(f(n))$ means ``grows at least as fast as $f(n)$'' --- a lower bound, a
        floor.
\end{itemize}
So we may also say: insertion sort runs in $O(n^2)$ on \emph{every} input (it is never
worse than quadratic), and in $\Omega(n)$ on every input (it must at least look at each
element once).

\begin{complexitybox}
\textbf{Insertion sort.}
\begin{itemize}
  \item Best case (sorted input): $\Theta(n)$ time.
  \item Worst case (reverse-sorted): $\Theta(n^2)$ time.
  \item Average case (random input): $\Theta(n^2)$ time.
  \item Extra space: $\Theta(1)$ --- it sorts \emph{in place}, using only a couple of
        variables beyond the array itself.
  \item Stable: yes (equal keys keep their original relative order).
\end{itemize}
\end{complexitybox}

\begin{aestheticbox}
You might think a quadratic sort is a museum piece. It is not. Insertion sort is the
fastest known method for \emph{small} arrays --- the constant factor $a$ is tiny because
each step is so cheap --- so real industrial sorting libraries (the ones inside your
phone, your browser, your database) use a fancy $\Theta(n\log n)$ method for big inputs
but \emph{switch to insertion sort} once the subproblems shrink below a threshold of
maybe $10$ to $30$ elements. It is also the natural choice when data arrives a little at
a time and is ``almost sorted'' already, since then it runs in nearly linear time. The
humble card-player's trick is quietly inside almost every program you use.
\end{aestheticbox}

\section{Designing algorithms: two stances}

Insertion sort is an example of the \emph{incremental} design stance: having sorted
$A[1\,..\,j-1]$, we insert one more element to extend the solution by a single step. We
grow the answer outward, one item at a time. Many algorithms fit this mold, and the
loop-invariant method is tailor-made for analyzing them.

But there is a second, often more powerful stance: \emph{divide and conquer}. Instead
of nibbling at the problem one element at a time, we split it into a few smaller
subproblems of the same type, solve each (often recursively), and then stitch the
sub-answers together. This stance frequently turns a $\Theta(n^2)$ algorithm into a
$\Theta(n \log n)$ one --- a colossal improvement at scale. Let us preview it with the
crown jewel of this chapter's successor: \emph{merge sort}.

\section{A first taste of merge sort}

Merge sort embodies divide and conquer in three crisp steps.

\begin{intuitionbox}
\textbf{Divide:} split the $n$-element array into two halves of about $n/2$ each.

\textbf{Conquer:} sort each half --- by calling merge sort on it recursively. A
one-element array is already sorted, so the recursion bottoms out there.

\textbf{Combine:} \emph{merge} the two sorted halves into one sorted whole.
\end{intuitionbox}

The clever, load-bearing piece is the \emph{merge}. Merging two already-sorted lists is
easy and fast: look at the front of each list, take the smaller of the two, and repeat.
Each comparison commits one element to the output, so merging two lists of total length
$n$ costs only about $n$ comparisons --- it is $\Theta(n)$, linear.

\begin{labbox}
\textbf{Merging $\langle 2,4,5,6\rangle$ and $\langle 1,3\rangle$.} Two fingers, one on
each list's front. Smaller finger's element goes to the output; advance that finger.

\smallskip
\begin{tabular}{l l l}
fronts $(1,2)$ & take $1$ & out: $1$ \\
fronts $(3,2)$ & take $2$ & out: $1,2$ \\
fronts $(3,4)$ & take $3$ & out: $1,2,3$ \\
right list empty & take rest $4,5,6$ & out: $1,2,3,4,5,6$ \\
\end{tabular}

\smallskip
Six output elements, and we never looked back. That ``one comparison per output
element'' is why merging is linear.
\end{labbox}

Figure~\ref{ch02:fig-merge} draws the same merge as two fingers walking down their
lists. At each step the smaller of the two pointed-at values is emitted and only that
finger advances; when one list runs dry, the rest of the other is copied wholesale.

\begin{figure}[H]
\centering
\begin{tikzpicture}[
    cell/.style={draw, minimum size=7mm, font=\small},
    srccell/.style={cell, fill=sourceblue!12},
    outcell/.style={cell, fill=targetgreen!16},
    takencell/.style={cell, fill=gray!12},
]
% left list
\node[font=\footnotesize,anchor=east] at (-0.2,1.1) {left:};
\foreach \v/\x/\sty in {2/0/takencell, 4/1/srccell, 5/2/srccell, 6/3/srccell}
  \node[\sty] at (\x*0.8,1.1) {\v};
% right list
\node[font=\footnotesize,anchor=east] at (-0.2,0.3) {right:};
\foreach \v/\x/\sty in {1/0/takencell, 3/1/takencell}
  \node[\sty] at (\x*0.8,0.3) {\v};
% fingers (current fronts after taking 1,2,3): left front = 4, right list exhausted
\draw[->,thick,quotientorange!75!black] (0.8,1.55) -- (0.8,1.35);
\node[font=\footnotesize,quotientorange!70!black,anchor=south] at (0.8,1.55) {finger};
\draw[->,thick,quotientorange!75!black] (1.6,-0.15) -- (1.6,0.05);
\node[font=\footnotesize,quotientorange!70!black,anchor=north] at (1.6,-0.15) {finger past end};
% output
\node[font=\footnotesize,anchor=east] at (3.6,-1.0) {output:};
\foreach \v/\x in {1/0,2/1,3/2,4/3,5/4,6/5}
  \node[outcell] at (3.6+\x*0.8,-1.0) {\v};
\end{tikzpicture}
\caption{Merging $\langle 2,4,5,6\rangle$ with $\langle 1,3\rangle$. Greyed cells were
already emitted; each comparison of the two finger cells sends one element to the output
below. Once the right list is exhausted, its finger walks off the end and the remaining
left cells $4,5,6$ are appended directly.}
\label{ch02:fig-merge}
\end{figure}

Now the cost. Let $T(n)$ be merge sort's running time on $n$ elements. Dividing is
cheap, $\Theta(1)$. Conquering is two recursive calls on inputs of size $n/2$, costing
$2\,T(n/2)$. Combining (the merge) is $\Theta(n)$. Writing $cn$ for the divide-plus-merge
work, we get the \emph{recurrence}
\[
T(n) = 2\,T(n/2) + cn, \qquad T(1) = \Theta(1).
\]
We will solve recurrences properly in the chapter on divide and conquer, but the
recursion tree of Figure~\ref{ch02:fig-tree} already gives the answer. Picture the
recursion as a tree: the root does $cn$ work splitting and merging the whole array; its
two children each do $c\,(n/2)$ work; their four children each do $c\,(n/4)$; and so on.

\begin{figure}[H]
\centering
\begin{tikzpicture}[
    level distance=11mm,
    level 1/.style={sibling distance=42mm},
    level 2/.style={sibling distance=21mm},
    every node/.style={font=\small},
    lvl/.style={draw, rounded corners, fill=sourceblue!10, inner sep=2.5pt},
]
\node[lvl] (r) {$cn$}
  child {node[lvl] {$c\,\tfrac{n}{2}$}
    child {node[lvl] {$c\,\tfrac{n}{4}$} child {node {$\vdots$}} child {node {$\vdots$}}}
    child {node[lvl] {$c\,\tfrac{n}{4}$} child {node {$\vdots$}} child {node {$\vdots$}}}}
  child {node[lvl] {$c\,\tfrac{n}{2}$}
    child {node[lvl] {$c\,\tfrac{n}{4}$} child {node {$\vdots$}} child {node {$\vdots$}}}
    child {node[lvl] {$c\,\tfrac{n}{4}$} child {node {$\vdots$}} child {node {$\vdots$}}}};
\node[anchor=west,gray] at (3.4,0) {\small level total: $cn$};
\node[anchor=west,gray] at (3.4,-1.1) {\small level total: $cn$};
\node[anchor=west,gray] at (3.4,-2.2) {\small level total: $cn$};
\end{tikzpicture}
\caption{The merge-sort recursion tree for $T(n)=2T(n/2)+cn$. The subproblem sizes on
each level sum back to $n$, so every level costs $cn$. With about $\log_2 n$ levels, the
total is $cn\log_2 n = \Theta(n\log n)$.}
\label{ch02:fig-tree}
\end{figure}

Here is the magic. \emph{Each level of the tree does $cn$ work in total}, because the
subproblem sizes at any level add back up to $n$ (the halves recombine). And there are
about $\log_2 n$ levels, because you can only halve $n$ down to $1$ about $\log_2 n$
times. Multiply: $(\text{work per level}) \times (\text{number of levels}) = cn \cdot
\log_2 n$, giving
\[
T(n) = \Theta(n \log n).
\]

\begin{focusbox}
Compare the two algorithms on $n = 1{,}000{,}000$. Insertion sort's worst case does
about $n^2 = 10^{12}$ operations. Merge sort does about $n \log_2 n \approx 10^6 \times 20
= 2 \times 10^7$. That is a factor of roughly \emph{fifty thousand}. No constant factor,
no faster laptop, ever closes a gap like that. This single comparison is why ``order of
growth'' is the first idea we taught --- it is the difference between an algorithm that
finishes and one that does not.
\end{focusbox}

\begin{historybox}
Merge sort dates to John von Neumann around 1945, making it one of the very first
algorithms designed for a stored-program computer. We will develop it in loving detail
in the chapter on divide and conquer, including a careful merge with sentinels, a full
correctness proof by loop invariant, and a rigorous solution of the recurrence $T(n) =
2T(n/2) + \Theta(n)$. Here it is only a teaser --- a promise that the quadratic wall is
not the end of the story. In CLRS this material lives in Section~2.3.
\end{historybox}

\begin{walkbox}
On your walk, weigh the trade. Merge sort is asymptotically far faster, but it needs
$\Theta(n)$ \emph{extra} memory to hold the merged output, whereas insertion sort sorts
in place with $\Theta(1)$ extra space. And for tiny arrays, insertion sort's small
constant actually wins. So ``which sort is best?'' has no single answer --- it depends on
$n$, on available memory, and on whether the data is nearly sorted already. When is each
one the right tool? Good algorithm design is exactly this kind of judgment.
\end{walkbox}

\section{What to carry forward}

\begin{recoverybox}
Before the exercises, gather the threads.
\begin{itemize}
  \item \textbf{The sorting problem} demands two things at once: a sorted output that is
        also a permutation of the input.
  \item \textbf{Insertion sort} grows a sorted region one element at a time, sliding each
        new key leftward into place --- the incremental design stance.
  \item \textbf{Loop invariants} prove correctness via initialization, maintenance, and
        termination. It is induction, plus a deliberate payoff at the end.
  \item \textbf{Running time} is counted in primitive steps as a function of input size
        $n$; insertion sort is $\Theta(n)$ at best, $\Theta(n^2)$ on average and at worst.
  \item \textbf{Order of growth} discards constants and lower-order terms to expose how
        cost scales --- the one feature that survives changing machines.
  \item \textbf{Divide and conquer}, previewed through merge sort, breaks $\Theta(n^2)$
        down to $\Theta(n \log n)$, a difference that dwarfs any constant factor.
\end{itemize}
\end{recoverybox}

\section{Exercises}

\begin{exercise}[Warm-up trace]
Trace \textsc{Insertion-Sort} on $A = \langle 31, 41, 59, 26, 41, 58 \rangle$, showing
the array after each pass of the outer loop. There are two $41$s; confirm from your
trace that the original (earlier) $41$ stays before the later one, illustrating
stability.
\end{exercise}

\begin{exercise}[Counting shoves]
For the array in the previous exercise, write down $t_j$ (the number of \textbf{while}
tests) for each $j$ from $2$ to $6$, and use your table to state the total number of
element shifts performed.
\end{exercise}

\begin{exercise}[Best and worst by construction]
Give a $6$-element array on which insertion sort performs the \emph{fewest} possible
shifts, and another on which it performs the \emph{most}. State the shift count in each
case, and connect them to the $\Theta(n)$ and $\Theta(n^2)$ analyses.
\end{exercise}

\begin{exercise}[Descending order]
Modify the pseudocode of \textsc{Insertion-Sort} so that it sorts into
\emph{non-increasing} order ($A[1] \ge A[2] \ge \cdots \ge A[n]$). Change as little as
possible, and state precisely which line(s) you changed and why.
\end{exercise}

\begin{exercise}[Linear search invariant]
Consider searching an array $A[1\,..\,n]$ for a value $v$, scanning left to right and
returning the index of the first match (or \textsc{Nil} if absent). Write the
pseudocode, then state a loop invariant and use it to argue correctness via
initialization, maintenance, and termination.
\end{exercise}

\begin{exercise}[The inner-loop invariant]
The maintenance step in this chapter waved its hand at the \textbf{while} loop. State a
precise loop invariant for the inner \textbf{while} loop of \textsc{Insertion-Sort} and
verify its initialization and maintenance. (Hint: describe what is true about
$A[i+1\,..\,j]$ and about the position of the gap.)
\end{exercise}

\begin{exercise}[Why the strict inequality]
Suppose someone changes the while-test from $A[i] > \mathit{key}$ to $A[i] \ge
\mathit{key}$. Argue that the algorithm is still correct (sorted $+$ permutation), but
exhibit a concrete input on which it does strictly more work, and explain why it is no
longer stable.
\end{exercise}

\begin{exercise}[Sum it carefully]
Prove by induction that $\sum_{k=1}^{n} k = \frac{n(n+1)}{2}$. Then use it to give the
\emph{exact} number of \textbf{while}-loop tests insertion sort performs on a
strictly-decreasing array of length $n$ (not just the order of growth).
\end{exercise}

\begin{exercise}[Binary insertion]
In the inner loop we find the insertion point by scanning left one step at a time. Since
$A[1\,..\,j-1]$ is already sorted, we could instead find the insertion point by
\emph{binary search} in $\Theta(\log j)$ comparisons. Explain why this reduces the
number of \emph{comparisons} to $\Theta(n \log n)$ but does \emph{not} reduce the
overall running time below $\Theta(n^2)$. (Hint: think about the cost of the shifting,
not the comparing.)
\end{exercise}

\begin{exercise}[Merge by hand]
Merge the sorted lists $\langle 1, 4, 7, 9 \rangle$ and $\langle 2, 3, 5, 8, 10 \rangle$
using the two-finger method, writing out the output after each step. Count the
comparisons and confirm the count is at most (length of first $+$ length of second)
$- 1$.
\end{exercise}

\begin{exercise}[Recursion-tree counting]
Merge sort on $n = 8$ elements. Draw the recursion tree, label the work at each node
($cn, c\frac{n}{2}, \dots$), and verify directly that the total over all nodes equals
$c \cdot 8 \cdot \log_2 8$. How many levels does the tree have?
\end{exercise}

\begin{exercise}[Crossover point, a challenge]
Suppose insertion sort costs exactly $\frac{1}{2} n^2$ ``time units'' in the worst case
and merge sort costs exactly $50\, n \log_2 n$ time units (merge sort's larger constant
reflects its overhead). For which values of $n$ is insertion sort actually \emph{faster}?
Solve the inequality $\frac{1}{2} n^2 < 50\, n \log_2 n$ approximately, and explain how
this justifies the ``switch to insertion sort for small subarrays'' trick used by real
libraries.
\end{exercise}
